\documentclass[12pt]{article}
\usepackage[utf8]{inputenc}
\usepackage{amsmath}
\usepackage{graphicx}
\usepackage{natbib}
\usepackage{hyperref}
\usepackage{lipsum} % For dummy text, you can remove this in your final document

\title{Effects of Noise Pollution on Marine Mammals: Physiological Stress Responses, Behavioral Changes, and Conservation Implications}
\author{Your Name}
\date{\today}

\begin{document}

\maketitle

\begin{abstract}
Noise pollution from anthropogenic sources poses significant threats to marine mammals, impacting their physiology and behavior. This paper reviews the physiological stress responses and behavioral changes observed in marine mammals due to noise pollution, analyzes case studies from field research and experimental studies, discusses conservation strategies, and provides recommendations for mitigating the effects of noise pollution on marine mammal populations.
\end{abstract}

\section{Introduction}
Marine mammals rely heavily on sound for communication, navigation, and foraging. Anthropogenic noise pollution from sources such as shipping, naval activities, and offshore development disrupts these vital functions, leading to adverse effects on marine mammal populations worldwide. This paper explores the impact of noise pollution on marine mammal physiology and behavior, discusses case studies highlighting specific effects, and proposes conservation strategies to mitigate these impacts.

\section{Impact Analysis}
\subsection{Physiological Stress Responses}
Noise pollution triggers physiological stress responses in marine mammals, including increased cortisol levels, altered reproductive hormones, and compromised immune function \citep{nowacek2007acoustic, rolland2012evidence}.

\subsection{Behavioral Changes}
Noise disturbance alters marine mammal behavior, affecting vocalization patterns, feeding habits, migration routes, and social interactions \citep{hatch2008evidence, wright2016marine}.

\section{Case Studies}
\subsection{Effects on Cetaceans}
Studies on cetaceans show avoidance behaviors, reduced foraging efficiency, and increased stress hormone levels in response to underwater noise from vessels \citep{southall2007marine, tyack2011beaked}.

\subsection{Effects on Pinnipeds}
Pinnipeds exhibit disrupted breeding behavior, altered haul-out patterns, and reduced pup survival rates due to noise pollution in their habitats \citep{acevedo2018effects, eguchi2019impact}.

\section{Conservation Strategies}
\subsection{Regulatory Measures}
Implementing and enforcing noise pollution regulations, including speed limits for vessels, quieting technologies, and designated quiet zones in critical habitats \citep{mccormick2019regulatory, bain2016mitigating}.

\subsection{Habitat Management}
Identifying and protecting important marine mammal habitats, minimizing human activities in sensitive areas, and promoting sustainable marine spatial planning \citep{costa2019habitat, smith2014habitat}.

\section{Recommendations}
Effective monitoring and research programs to assess noise impacts, public awareness campaigns on the importance of reducing noise pollution, and international collaboration for global conservation efforts \citep{simmonds2018marine, sparrow2020global}.

\section*{References}
\bibliographystyle{plainnat}
\bibliography{prompt10_35}

\end{document}
